\section*{Capítulo \ref{ch:b3:radicals-carbenes} --- Radicales, carbenos y transposiciones}

\begin{solution}{exo:b3:radicals-carbenes:1}
El propano tiene seis hidrógenos primarios y dos secundarios. Cloración: $6 \times 1 : 2 \times 3.9 = 6 : 7.8$,
es decir, un 43\,\% de 1-cloropropano y un 57\,\% de 2-cloropropano. Bromación: $6 : 164$, un 4\,\% y un
96\,\%.
\end{solution}

\begin{solution}{exo:b3:radicals-carbenes:2}
$\ce{CH2}$: triplete; $\ce{CCl2}$: singlete (los pares solitarios del cloro elevan el orbital p vacío).
El diclorocarbeno se adiciona de forma estereoespecífica: \textit{trans}-1,1-dicloro-2,3-dimetilciclopropano.
\end{solution}

\begin{solution}{exo:b3:radicals-carbenes:3}
\textit{cis}-1,2-Dietilciclopropano: el \omterm{def:b3:radicals-carbenes:carbenoid}{carbenoide} se adiciona en una sola etapa por una sola cara.
\end{solution}

\begin{solution}{exo:b3:radicals-carbenes:4}
Ciclohexanona: $\varepsilon$-caprolactona (oxepan-2-ona). Acetofenona: acetato de fenilo,
$\ce{CH3CO2C6H5}$; el fenilo se desplaza con preferencia al metilo.
\end{solution}

\begin{solution}{exo:b3:radicals-carbenes:5}
Nueve hidrógenos primarios y uno terciario. Cloración: $9 : 5.1$, un 64\,\% de 1-cloro-2-metilpropano
y un 36\,\% de 2-cloro-2-metilpropano. Bromación: $9 : 1600$, un 0.6\,\% y un 99.4\,\%.
\end{solution}

\begin{solution}{exo:b3:radicals-carbenes:6}
$k_c = 3.0 \times k_H[\mathrm{Bu_3SnH}] = 3.0 \times \num{2.4e6} \times 0.020 = \qty{1.4e5}{s^{-1}}$.
\end{solution}

\begin{solution}{exo:b3:radicals-carbenes:7}
Ataque sobre C5: 5-exo-trig, favorecido; sobre C6: 6-endo-trig, favorecido por las reglas pero
más lento aquí. 4-exo-tet: favorecido. 5-endo-trig: desfavorecido. 5-exo-dig: favorecido.
3-exo-dig: desfavorecido. 5-endo-dig: favorecido.
\end{solution}

\begin{solution}{exo:b3:radicals-carbenes:8}
La pérdida de agua da un catión terciario y bencílico; en el carbono vecino podrían desplazarse un fenilo
y un grupo metilo, y se desplaza el fenilo, de mayor aptitud. El
catión estabilizado por el oxígeno pierde un protón: 3,3-difenilbutan-2-ona.
\end{solution}

\begin{solution}{exo:b3:radicals-carbenes:9}
Curtius: $\ce{C6H5CON3 -> C6H5NCO + N2}$. Con agua, el ácido carbámico pierde $\ce{CO2}$: anilina
(se forma algo de difenilurea cuando la anilina encuentra más \omterm{def:b3:radicals-carbenes:curtius}{isocianato}). Con etanol: \textit{N}-fenilcarbamato
de etilo, $\ce{C6H5NHCO2C2H5}$.
\end{solution}

\begin{solution}{exo:b3:radicals-carbenes:10}
Cloro: primario, $421.7 - 432 = \qty{-10}{kJ/mol}$; terciario, $405.4 - 432 = \qty{-27}{kJ/mol}$; bromo: $+56$ y
$+39$. Las dos abstracciones por el cloro son exotérmicas, con estados de transición tempranos que
apenas perciben la diferencia de \qty{16}{kJ/mol}; las dos abstracciones por el bromo son endotérmicas,
con estados de transición tardíos cuyas energías de activación difieren en la mayor parte de ella:
$\eu^{16000/(8.314 \times 298)} \approx 600$, el orden del cociente observado.
\end{solution}

\begin{solution}{exo:b3:radicals-carbenes:11}
Oxima \textit{E} (OH en anti respecto de C2, que lleva el grupo metilo): se desplaza C2, y el nitrógeno entra
junto a él: 7-metilazepan-2-ona. Oxima \textit{Z}: se desplaza C6: 3-metilazepan-2-ona. En la
\omterm{def:b3:radicals-carbenes:beckmann}{transposición de Beckmann} decide la geometría de la oxima; en la \omterm{def:b3:radicals-carbenes:baeyer}{oxidación de Baeyer--Villiger},
la \omterm{def:b3:radicals-carbenes:migratory}{aptitud migratoria}, y se desplaza el carbono secundario:
7-metiloxepan-2-ona.
\end{solution}

\begin{solution}{exo:b3:radicals-carbenes:12}
Fracción ciclada: $k_c/(k_c + k_H c)$: 0.91, 0.49 y 0.09. Para el 95\,\%: $k_Hc = k_c(1/0.95 - 1)$, luego
$c = 0.0526 \times \num{2.3e5}/\num{2.4e6} = \qty{5.0e-3}{mol/L}$. Se mantiene añadiendo el hidruro lentamente con una
bomba de jeringa, o usando una cantidad catalítica de haluro de estaño con un reductor
que regenera el hidruro.
\end{solution}

\begin{solution}{pb:b3:radicals-carbenes:1}
\textbf{1.} $\ce{C6H10O + NH2OH -> C6H11NO + H2O}$.

\textbf{2.} Hace falta ácido para eliminar en forma de agua el grupo hidroxi del intermedio de adición,
pero demasiado ácido protona la hidroxilamina, que entonces ya no se adiciona.

\textbf{3.} No: los dos carbonos unidos al carbono C=N son equivalentes por la
simetría del anillo.

\textbf{4.} En la oxima \textit{E}, el OH apunta en sentido contrario al carbono C2, que lleva el metilo (en anti respecto de
él); en la oxima \textit{Z}, hacia él.

\textbf{5.} La inversión en un nitrógeno de oxima que lleva un oxígeno electronegativo tiene una barrera
alta.

\textbf{6.} Protona (o sulfona) el grupo hidroxi y lo convierte en un buen grupo
saliente, el agua.

\textbf{7.} El enlace C--C del anillo situado en anti respecto del grupo saliente pasa del carbono al nitrógeno,
que queda así insertado en el anillo: seis átomos pasan a ser siete.

\textbf{8.} Un ion nitrilio cíclico.

\textbf{9.} $\ce{C6H11NO -> C6H11NO}$: una isomerización, de oxima a lactama.

\textbf{10.} Caprolactama, azepan-2-ona, la amida cíclica de siete miembros del
ácido 6-aminohexanoico; su polimerización por apertura de anillo da el nailon 6.

\textbf{11.} \qty{98.0}{g/mol} y \qty{113.0}{g/mol}.

\textbf{12.} $\qty{1.0e6}{g}/\qty{98.0}{g/mol} = \qty{10204}{mol}$; $\times 0.98 \times 0.98 = \qty{9800}{mol}$; $\times \qty{113.0}{g/mol} = \qty{1.11e6}{g}$, \qty{1.11}{t}.

\textbf{13.} Oxima: $\num{10204} \times 0.98 = \qty{1.00e4}{mol}$; $\ce{H2SO4}$: $\qty{1.30e4}{mol}$, que dan otros tantos moles de
$\ce{(NH4)2SO4}$: $\qty{1.30e4}{mol} \times \qty{132.1}{g/mol} = \qty{1.7}{t}$, más que la caprolactama.

\textbf{14.} $\ce{C6H10O + RCO3H -> C6H10O2 + RCOOH}$.

\textbf{15.} El aducto tetraédrico del peroxiácido sobre el carbono carbonílico (el
intermedio de Criegee), un peroxihemicetal.

\textbf{16.} Un grupo $\ce{CH2}$ del anillo (los dos son equivalentes) se desplaza hacia el oxígeno peroxídico
más cercano; sale el carboxilato $\ce{RCOO-}$.

\textbf{17.} $\varepsilon$-Caprolactona, oxepan-2-ona; su polimerización por apertura de anillo da la
poli($\varepsilon$-caprolactona), un poliéster biodegradable.

\textbf{18.} Su único subproducto es el agua, mientras que un peroxiácido deja un mol de
ácido carboxílico por mol de lactona y es en sí mismo peligroso de almacenar.

\textbf{19.} Se desplaza el carbono secundario C2: 7-metiloxepan-2-ona.

\textbf{20.} Sí: el carbono que se desplaza conserva su configuración.

\textbf{21.} 7-Metilazepan-2-ona (C2 en anti, se desplaza).

\textbf{22.} 3-Metilazepan-2-ona (se desplaza C6).

\textbf{23.} Beckmann: la geometría de la oxima (se desplaza el grupo anti).
Baeyer--Villiger: la \omterm{def:b3:radicals-carbenes:migratory}{aptitud migratoria} (se desplaza el carbono más sustituido).

\textbf{24.} Sus dos carbonos $\alpha$ son equivalentes: una sola oxima, una sola lactama, una sola lactona.

\textbf{25.} Una tonelada de ciclohexanona da unas \qty{1.11}{t} de caprolactama.
\end{solution}
